\documentclass{tufte-handout}

%\geometry{showframe}% for debugging purposes -- displays the margins



\newcommand{\bra}[1]{\left(#1\right)}
\usepackage{clrscode3e}
\usepackage{tikz}
\usepackage{amsmath,amsthm}
\usetikzlibrary{shapes}
\usetikzlibrary{positioning}
% Set up the images/graphics package
\usepackage{graphicx}
\setkeys{Gin}{width=\linewidth,totalheight=\textheight,keepaspectratio}
\graphicspath{{graphics/}}

\title{Directed Reading Project writeup}
\author{Shan Gao}
\date{Winter 2024}  % if the \date{} command is left out, the current date will be used

% The following package makes prettier tables.  We're all about the bling!
\usepackage{booktabs}

% The units package provides nice, non-stacked fractions and better spacing
% for units.
\usepackage{units}

% The fancyvrb package lets us customize the formatting of verbatim
% environments.  We use a slightly smaller font.
\usepackage{fancyvrb}
\fvset{fontsize=\normalsize}

% Small sections of multiple columns
\usepackage{multicol}

%--------Theorem Environments--------
%theoremstyle{plain} --- default
\newtheorem{thm}{Theorem}
\newtheorem{cor}[thm]{Corollary}
\newtheorem{prop}[thm]{Proposition}
\newtheorem{lem}[thm]{Lemma}
\newtheorem{conj}[thm]{Conjecture}
\newtheorem{quest}[thm]{Question}

\theoremstyle{definition}
\newtheorem{defn}[thm]{Definition}
\newtheorem{defns}[thm]{Definitions}
\newtheorem{con}[thm]{Construction}
\newtheorem{exmp}[thm]{Example}
\newtheorem{exmps}[thm]{Examples}
\newtheorem{notn}[thm]{Notation}
\newtheorem{notns}[thm]{Notations}
\newtheorem{addm}[thm]{Addendum}
\newtheorem{exer}[thm]{Exercise}

\theoremstyle{remark}
\newtheorem{rem}[thm]{Remark}
\newtheorem{rems}[thm]{Remarks}
\newtheorem{warn}[thm]{Warning}
\newtheorem{sch}[thm]{Scholium}


% These commands are used to pretty-print LaTeX commands
\newcommand{\doccmd}[1]{\texttt{\textbackslash#1}}% command name -- adds backslash automatically
\newcommand{\docopt}[1]{\ensuremath{\langle}\textrm{\textit{#1}}\ensuremath{\rangle}}% optional command argument
\newcommand{\docarg}[1]{\textrm{\textit{#1}}}% (required) command argument
\newenvironment{docspec}{\begin{quote}\noindent}{\end{quote}}% command specification environment
\newcommand{\docenv}[1]{\textsf{#1}}% environment name
\newcommand{\docpkg}[1]{\texttt{#1}}% package name
\newcommand{\doccls}[1]{\texttt{#1}}% document class name
\newcommand{\docclsopt}[1]{\texttt{#1}}% document class option name

\begin{document}
\maketitle
\section*{Preliminaries}
\subsection*{Basic definitions}
\begin{defn}[Algebra]
Let $S$ be a set, an algebra on $S$ , notated $\Sigma_0$, is a collection of subsets of $S$ that satisfies the following conditions:
\begin{itemize}
    \item $S \in \Sigma_0$
    \item $F \in \Sigma_0 \implies F^c :=  S \setminus F \in \Sigma_0$
    \item $F,G \in \Sigma_0 \implies F \cup G \in \Sigma_0$
    
\end{itemize}
\end{defn}
\begin{defn}[$\sigma$-algebra]
A collection $\Sigma$ of subsets of $S$ is an algbra on $S$ if $\Sigma$ is an algebra on $S$ and whenever $\{F_n\}_{n \in \mathbb{N}} \subseteq \Sigma$ then
$$
\bigcup_{n} F_n \in \Sigma
$$
So it's really just an algebra with the added condition that it's closed under \textbf{countable} union

    
\end{defn}
\begin{defn}[Measurable space]
    A pair $(S,\Sigma)$, where $S$ is a set and $\Sigma$ is a $\sigma$-algebra on $S$, is called a measurable space. An element of $\Sigma$ is called a $\Sigma$-measurable subset of $S$. 
\end{defn}
\begin{defn}[algbera generated by class of subsets]
    Let $\mathcal{C}$ be a class of subsets of $S$. Then $\sigma(\mathcal{C}))$, the $\sigma$-algebra generated by $\mathcal{C}$, is the smallest $\sigma$-algebra $\Sigma$ on $S$ such that $\mathcal{C} \subseteq \Sigma$. Precisely, it is the intersection of all $\sigma$-algberas on $S$ which have $\mathcal{C}$ as a subclass. 
\end{defn}
\begin{exmp}[Borel $\sigma$-algebras]
We call \textbf{Borel} $\sigma$-algebra on $S$, denoted $\mathcal{B}(S)$, to tbe the $\sigma$-algebra generated by the collection of all open subsets of $S$. 
$$
\mathcal{B}(S) := \sigma(\text{open sets of $S$} )
$$
We call the Borel $\sigma$-algebra the $\sigma$-algebra generated by all open subsets of the reals, where it is standard to notate $\mathcal{B} = \mathcal{B}(\mathbb{R})$
    
\end{exmp}
\begin{defn}[Additive, countably additive]
    Let $S$ be a set, and let $\Sigma_0$ be an algebra on $S$, and let $\mu_0$ be a non-negative set function 
    $$
        \mu_0 \colon \Sigma_0 \to [0,\infty]
    $$
    
    $\mu_0$ is additive if $\mu(\emptyset) = 0$ and, for $F,G \in \Sigma_0$, 

    $$
    F \cap G = \emptyset \implies \mu_0 (F \cup G) = \mu_0(F) + \mu_0(G)
    $$
    it's \textbf{countably additive} or $\sigma$-additive if $\mu(\emptyset) = 0 $ and whenever $(F_n \colon n \in \mathbb{N})$ is a sequence of disjoint sets in $\Sigma_0$ with union $F = \bigcup_{n} F_n$ in $\Sigma_0$, then
    $$ \mu_0(F) = \sum_{n} \mu_0(F)$$

    
\end{defn}

\begin{defn}[Measure space]
Let $(S, \Sigma)$ be a measurable space, so that $\Sigma$ is a $\sigma$-algebra on $S$. A map 
$$\mu \colon \Sigma \to [0,\infty)$$
that is countably additive forms a measure on $(S,\Sigma)$. The triple $(S, \Sigma, \mu)$ is then called a measure space. 
    
\end{defn}
\begin{defn}[Probability measure, probability triple]
    Our measure $\mu$ is called a probability measure if 
    $$\mu(S) = 1$$
    and then $(S, \Sigma, \mu)$ is then called a probability triple.
\end{defn}
\begin{defn}[$\sigma$-null element of $\Sigma$, almost everywhere (a.e)]
	An element $F$ of $\Sigma$ is called $\mu$-null if $\mu(F) = 0$. A statement $\mathcal{S}$ about points $s$ of $S$ is said to hold almost everywhere (a.e) if
	$$
		F := \{s \colon \mathcal{S}(s) \text{ is false }\} \in \Sigma \text { and } \mu(F) = 0
	$$
\end{defn}
\begin{lem}[ $\pi$-systems and uniqueness of extensions]
	Let $S$ be a set. Let $I$ be a $\pi$-system on $S$, that is, a family of subsets of $S$ closed under finite intersection,  then letting $\Sigma = \sigma(I)$. If $\mu_1$ and $\mu_2$ are measures on $(S, \Sigma)$ such that $\mu_1(S) = \mu_2(S) < \infty$ and $\mu_1 = \mu_2$ on $I$, then 
	$$\mu_1 = \mu_2$$ on $\Sigma$. In human terms, if two measures agree on a $\pi$-system then they agree on the $\sigma$-algebra generated by that $\pi$-system.
\end{lem}
\begin{thm}[Caratheodory's Extension Theorem]
	Let $S$ be a set, let $\Sigma_0$ be an algebra on $S$, and let 
	$$\Sigma := \sigma(\Sigma_0)$$
	If $\mu_0$ is a countably additive map $\mu_0 \colon \Sigma_0 \to [0,\infty]$, then there exists a measure $\mu$ on $(S,\Sigma)$ such that 
	$$\mu = \mu_0$$ on $\Sigma$, if the map is finite on $S$, then this extension is unique. In human terms, if we have an algebra and a countably additive map on the algebra, then there exists an extension to a measure space that agrees with the map. 
\end{thm}

\begin{lem}[Elementary inequalities]
	Let $(S, \Sigma, \mu)$ be a measure space. Then 
	\begin{enumerate}
		\item[(a)] $\mu (A \cup B) \leq \mu(A) + \mu(B)$ with $A,B \in \Sigma$
		\item[(b)] $\mu( \bigcup_{i \leq n} F_i) \leq \sum_{i \leq n} \mu(F_i)$ ($F_1, F_2, \ldots, F_n \in \Sigma$)
			Furthermore, if $\mu(S) < \infty$, then 
		\item[(c)] $\mu(A \cup B) = \mu(A) + \mu(B) - \mu(A \cap B)$ for $(A,B \in \Sigma)$
		\item[(d)] \textbf{inclusion exclusion} : for $F_1, F_2, \ldots, F_n \in \Sigma$, 
				$\mu\left(\bigcup_{i \leq n} F_i \right) = \sum_{i \leq n} \mu(F_i) - \sum \sum_{i < j \leq n} \mu(F_i \cap F_j)  +\sum\sum\sum_{i < j < k \leq n } \mu(F_i \cap F_j \cap F_k) - \cdots + (-1)^{n-1} \mu(F_1 \cap F_2 \cap \ldots \cap F_n)$
	
	\end{enumerate}
	
	
\end{lem}
\begin{lem}[Continuity of measure]
	if $F_n \in \sigma ( n \in \mathbb{N})$  and $F_n \uparrow F$, then $\mu(F_n) \uparrow \mu(F)$. Same with if $G_n \in \Sigma (n \in \mathbb{N})$ and $G_n \downarrow G$ and $\mu(G_k) < \infty$ for some $k$, then $\mu(G_n) \downarrow \mu(G)$
\end{lem}
We introduce some probability notions, a probability triple denoted $(\Omega, \mathcal{F}, \mathbb{P})$ can model an experiment invovling randomness, where $\Omega$ is the \textbf{sample space}, $\omega \in \Omega$ is called a \textbf{sample point} and an event $F \in \mathcal{F}$ is just an $\mathcal{F}$ measurable subset of the $\sigma$-algebra $\mathcal{F}$ called the \textbf{family of events}. 
\begin{defn}[Almost surely (a.s.)]
	A statement $S$ about outcomes is said to be true \textbf{almost surely (a.s.)}, or with probability $1$ (w.p.1), if 
	$$ F := \{ \omega \colon S(\omega) \text{ is true } \} \in \mathcal{F} \text{ and } \mathbb{P}(F) = 1$$
	
\end{defn}
\begin{defn}[$\limsup,\liminf$ and more magic]
	Let $(x_n \colon n \in \mathbb{N})$ be a sequence of real numbers. We define 
	$$\limsup x_n := \inf_{m} \{\sup_{n \geq m} x_n \} = \downarrow \lim_{m} \{ \sup_{n \geq m } x_n \} \in [-\infty,+\infty]$$
	Obviously, $y_m := \sup_{n \geq m } x_n$ is monotone non-increasing in $m$, so that the limit of the sequence $y_m$ exists in $[-\infty, \infty]$. $\uparrow$ and $\downarrow$ are used to signify monotone limits, where if you have a sequence of numbers $(y_n \colon n \in \mathbb{N})$, then $y_n \downarrow y_{\infty}$ means the monotone limit $y_n \downarrow \lim_{n} y_n$.
\end{defn}
\begin{defn}[upper lattice]
	This is called the $\limsup_{n} E_n$ of some sequence of sets $(E_n \in \mathcal{F} \colon n \in \mathbb{N})$ to be a sequence of $\mathcal{F}$-measurable subsets of $\Omega$, then 
	\begin{align*}
		\mathcal{F} \ni (E_n \text{ infinitely often } )  &= \limsup_{n} E_n \\
								  &= \bigcap_{m} \bigcup_{n \geq m} E_n\\
								  &= \{ \omega \colon \forall m, \exists n(\omega) \geq m \text{ s.t } \omega \in E_{n(\omega)}\}\\
								  &= \{ \omega \colon \omega \in E_n \text{ for infinitely many $n$ } \}
	\end{align*}
	

\end{defn}
\begin{defn}[Reverse Fatou Lemma]
	This is true for any probability measure, since $\mathbb{P}$ is always $\sigma$-finite. 
	$$\mathbb{P}(\limsup E_n) \geq \limsup \mathbb{P}(E_n)$$
		
\end{defn}
\begin{lem}[First Borel-Canteli Lemma (BC1)]
	Let $(E_n \colon n \in \mathbb{N})$ be a sequence of events such that $\sum_{n} \mathbb{P}(E_n) < \infty$, then 
	$$ \mathbb{P}(\limsup_{n} E_n) = \mathbb{P}(E_n \text{ infinitely often }) = 0$$

\end{lem}

\end{document}
